\documentclass[UTF8]{ctexart}
\usepackage{amsmath}
\usepackage{amssymb}
\usepackage{geometry}
\usepackage{listings}
\geometry{a4paper, margin=1in}

\title{计算机图形学模拟试卷（三）参考答案与解析}
\date{}

\begin{document}

\maketitle

\section{简答题建议答案}
\begin{enumerate}
    \item \textbf{图形 vs 图像}：图形是基于算法生成的几何表示（抽象数据 $\to$ 图像）；图像处理是对已有像素点阵的加工（图像 $\to$ 图像）。二者在计算机视觉和现代渲染中相互渗透。
    \item \textbf{D3D12 流程}：CPU 准备资源 $\to$ 将资源（顶点等）上传到 GPU 堆 $\to$ 绑定根签名和 PSO $\to$ 发布记录命令列表 $\to$ 命令队列执行 $\to$ GPU 渲染。
    \item \textbf{顶点结构}：位置（Position，几何形状）、颜色（Color）、法线（Normal，光照计算）、切线（复杂着色）、纹理坐标（UV 映射）。
\end{enumerate}

\section{计算题解析}
\begin{enumerate}
    \item \textbf{Bresenham (10, 3)}：$\Delta x=10, \Delta y=3$。$d = 2\Delta y - \Delta x = 6 - 10 = -4$。
    1. (0,0), $d=-4 < 0 \to d = -4+6=2$。
    2. (1,0), $d=2 \ge 0 \to y=1, d = 2+2(3-10)=-12$。
    3. (2,1), $d=-12 < 0 \to d = -12+6=-6$ $\to$ (3,1)。
    
    \item \textbf{3D 变换}：$M = T(10, 20, 30) \cdot R(\theta)$。遵循“先变换、后叠加”顺序。
    
    \item \textbf{C-S 编码}：窗口 $[10,20]$。
    $P_1(5, 15)$：$x < 10 \implies$ 0001。
    $P_2(25, 25)$：$x > 20, y > 20 \implies$ 1010。
    不可简弃，相穿。
\end{enumerate}

\section{画图与分析解析}
\begin{itemize}
    \item \textbf{纹理寻址}：
    \begin{itemize}
        \item Wrap: $u \mod 1.0$ (贴瓷砖效果)。
        \item Mirror: 坐标在整数边界处翻转。
        \item Clamp: 强制为 0 或 1（取边缘色）。
    \end{itemize}
    \item \textbf{De Casteljau}：
    1. 在 $P_0P_1$ 线段中点找到 $M_{01}$。
    2. 在 $P_1P_2$ 线段中点找到 $M_{12}$。
    3. 在 $M_{01}M_{12}$ 线段中点找到 $M_{final}$（此即曲线上点）。
\end{itemize}

\section{算法解析}
\begin{lstlisting}[language=C++]
void UpdatePixelZ(int x, int y, float newZ, Color newColor) {
    if (newZ < ZBuffer[x][y]) {
        ZBuffer[x][y] = newZ;
        FrameBuffer[x][y] = newColor;
    }
}
\end{lstlisting}

\end{document}
